\section{Proofs in Section \ref{sec:DP-O2NC_privacy}}

\subsection{Proof of Corollary \ref{cor:DP-O2NC:privacy}}

% \TreeCorollary*

\begin{proof}
Since the private oracle defined in Algorithm \ref{alg:DP-O2NC:oracle-variance-reduced} uses disjoint data in different iteration $k\in[K]$, Algorithm \ref{alg:DP-O2NC:PONC} is a post-processing of $\{(\tilde\vecg_1^k,\ldots,\tilde\vecg_T^k)\}_{k\in[K]}$. Therefore, it suffices to prove that for each $k\in[K]$, the composition $(\tilde\vecg_1^k,\ldots,\tilde \vecg_T^k)$ is $(\alpha,\alpha\rho^2/2)$-RDP.

The key is to observe that $(\tilde\vecg_1^k,\ldots,\tilde\vecg_T^k)$ can be written as the adaptive composition in Theorem \ref{thm:DP-O2NC:tree}. Let $\vecg_1^k = \calM^{(1)}(Z_1^k) \triangleq \grad_{f,\delta}(\vecw_1^k,Z_1^k)$ and $\vecd_i^k = \calM^{(i)}(\tilde\vecg_{1:i-1}^k, Z_i^k) \triangleq \diff_{f,\delta}(\vecw_i^k,\vecw_{i-1}^k,Z_i^k)$. This is well-defined because 
% $\vecw_1^k$ is initialized by the OCO algorithm and is thus data-independent, and 
$\vecw_i^k$ is a post-processing of $\tilde\vecg_{1:i-1}^k$ by construction of O2NC. Therefore, $\vecg_t^k=\vecg_1^k+\sum_{i=2}^t \vecd_i^k = \sum_{i=1}^t \calM^{(i)}(\tilde\vecg_{1:i-1}^k,Z_i^k)$, and Theorem \ref{thm:DP-O2NC:tree} can be applied.

By Lemma \ref{lem:DP-O2NC:grad-estimator}, the sensitivity of $\calM^{(1)}$ is bounded by $s_1 = \frac{2dL}{B_1}$. By Remark \ref{rmk:DP-O2NC:w-distance} and Lemma \ref{lem:DP-O2NC:diff-estimator}, the sensitivity of $\calM^{(i)}(\tilde\vecg_{1:i-1}^k,\cdot)$ is bounded by $s_i=\frac{2dL}{B_2\delta}\|\vecw_i^k-\vecw_{i-1}^k\| \le \frac{4dL}{B_2T}$ for all $i\ge 2$.
% because $\|\vecw_i^k-\vecw_{i-1}^k\| \le 2D = \frac{2\delta}{T}$. 
Since we assume $B_1\ge TB_2/2$, it holds that $s_1\le s_t$ for all $t\ge 2$. Therefore, upon setting $\sigma_t = {\sqrt{2\ln T}s_2}/{\rho}$ for all $t\in[T]$, Theorem \ref{thm:DP-O2NC:tree} implies that $(\tilde\vecg_1^k,\ldots\tilde\vecg_T^k)$ is $(\alpha,\alpha\rho^{\prime 2}/2)$-RDP where
\begin{equation*}
    % \textstyle
    \rho' \le \sqrt{2\ln T} \cdot \max_{b\in[n], i\le n} \frac{s_i}{\sigma_b} \le \sqrt{2\ln T} \cdot \frac{s_2}{\sigma_2} = \rho.
    \qedhere
\end{equation*}
\end{proof}


\subsection{Proof of Corollary \ref{cor:DP-O2NC:tree-noise}}

% \TreeNoise*

\begin{proof}
Recall that $\tree(t)=\sum_{(\cdot,i)\in\node(t)} \xi_i$. Since $\xi_i$'s are independent and $\xi_i\sim \calN(0,\sigma_i^2I)$,
\begin{equation*}
    \E\|\tree(t)\|^2
    = \sum_{(\cdot,i)\in\node(t)} \E\|\xi_i\|^2
    = \sum_{(\cdot,i)\in\node(t)} d\sigma_i^2
    \le 2\ln(t) d\sigma^2.
\end{equation*}
The last inequality follows from Remark \ref{rmk:DP-O2NC:tree} that $|\node(t)|\le 2\ln(t)$. We then conclude the proof by substituting the value of $\sigma$ defined in Corollary \ref{cor:DP-O2NC:privacy}.
\end{proof}